% Festlegung des allgemeinen Dokumentenformats
\documentclass[a4paper,12pt]{article}

% Schrift
\usepackage[T1]{fontenc}
\usepackage{lmodern}
\usepackage[utf8]{inputenc}
\usepackage[ngerman]{babel}

% Bilder
\usepackage{graphicx}
\usepackage{float}
\graphicspath{{./assets/img}}

% Variablen
% Variablen
\newcommand{\titleDocument}{Wissenschaftliche Abhandlung}
\newcommand{\subjectDocument}{Realisierbarkeit eines Upload-Filters in Bezug auf das Urheberrechts-Diensteanbieter-Gesetz am Beispiel von Grafiken}
\newcommand*{\bildquelle}{
  \footnotesize Quelle:
}
\newcommand{\code}[1]{\noindent\ignorespaces\texttt{#1}}

% Code
\usepackage{listings}
\usepackage{color}

\definecolor{black}{rgb}{0,0,0}
\definecolor{dkgreen}{rgb}{0,0.6,0}
\definecolor{gray}{rgb}{0.5,0.5,0.5}
\definecolor{mauve}{rgb}{0.58,0,0.82}

\lstset{literate=%
    {Ö}{{\"O}}1
    {Ä}{{\"A}}1
    {Ü}{{\"U}}1
    {ß}{{\ss}}1
    {ü}{{\"u}}1
    {ä}{{\"a}}1
    {ö}{{\"o}}1
    {~}{{\textasciitilde}}1
}

\lstdefinestyle{Bash} {
  frame=single,
  language=Bash,
  aboveskip=3mm,
  belowskip=3mm,
  showstringspaces=false,
  columns=flexible,
  basicstyle={\small\ttfamily},
  numbers=none,
  numberstyle=\tiny\color{black},
  keywordstyle=\color{black},
  commentstyle=\color{black},
  stringstyle=\color{black},
  breaklines=true,
  breakatwhitespace=true,
  tabsize=3
}

\lstset{emph={\$},emphstyle=\textbf}

\lstdefinestyle{Python} {
  frame=single,
  language=Bash,
  aboveskip=3mm,
  belowskip=3mm,
  showstringspaces=false,
  columns=flexible,
  basicstyle={\small\ttfamily},
  numbers=left,
  numberstyle=\tiny\color{mauve},
  keywordstyle=\color{mauve},
  commentstyle=\color{dkgreen},
  stringstyle=\color{dkgreen},
  breaklines=true,
  breakatwhitespace=true,
  tabsize=2
}

% mehrseitige Tabellen ermöglichen
\usepackage{longtable}
\usepackage{diagbox}

% Packet für Seitenrandabstände und Einstellung für Seitenränder
\usepackage{geometry}
% Internet
%\geometry{left=3.5cm, right=2.5cm, top=2.5cm, bottom=2cm}
% Kern
\geometry{a4paper, top=27mm, left=20mm, right=20mm, bottom=35mm, headsep=10mm,
footskip=12mm}

% bricht lange URLs "schön" um
\usepackage[hyphens,obeyspaces,spaces]{url}

% Festlegung Art der Zitierung
\usepackage{csquotes}
\usepackage[style=apa, backend=biber]{biblatex}
\addbibresource{assets/literatur.bib}

% Abstand zwischen Absätze
\setlength{\parindent}{2em}
\setlength{\parskip}{1em}

% Paket für Zeilenabstand
\usepackage{setspace}
\onehalfspacing
% Kern:
% \setstretch{1.15}

% für Bildbezeichner
\usepackage{capt-of}

% für Stichwortverzeichnis
\usepackage{makeidx}

% für Abkürzungsverzeichnis
\usepackage{acronym}

% Für Phantomsection
\usepackage{hyperref}

% Für Tabellen
\usepackage{tabularx}

% Konfiguriere das Inhaltsverzeichnis
\usepackage{tocbasic}

% Anhang: PDF einfügen
\usepackage{pdfpages}

% subsubsubsection durch paragraph
\usepackage{titlesec}
\setcounter{secnumdepth}{4}
\titleformat{\paragraph}
{\normalfont\normalsize\bfseries}{\theparagraph}{1em}{}
\titlespacing*{\paragraph}
{0pt}{3.25ex plus 1ex minus .2ex}{1.5ex plus .2ex}
% Kern:
\titlespacing{\section}{0pt}{12pt plus 4pt minus 2pt}{8pt plus 2pt minus 2pt}
\titlespacing{\subsection}{0pt}{12pt plus 4pt minus 2pt}{6pt plus 2pt minus 2pt}
\titlespacing{\subsubsection}{0pt}{12pt plus 4pt minus 2pt}{4pt plus 2pt minus 2pt}

% \autoref: Subsubsection umbenennen
\addto\extrasngerman{
    \def\subsubsectionautorefname{Unterabschnitt}
}

% Titel
\title{Bachelorarbeit}

% Autor
\author{Andreas Huber}

% Datum
\date{\today}

%
% Start
% des
% Dokuments
%
\begin{document}

% Titelseite
\thispagestyle{empty}

\begin{figure}[t]
 \centering
 \includegraphics[width=0.4\textwidth]{assets/oth/logo}
\end{figure}

\begin{verbatim}
\end{verbatim}

\begin{center}
    \Large{Ostbayerische Technische Hochschule Regensburg} \\
    \Large{Fakultät für Informatik und Mathematik}
\end{center}

\begin{verbatim}
\end{verbatim}

\begin{center}
    \doublespacing
    \textbf{\huge{\titleDocument}}\\

    \onehalfspacing

    % \begin{center}
    %     Zur Erlangung des akademischen Grades \\ Bachelor of Science (B. Sc.)
    % \end{center}

    \begin{verbatim}
    \end{verbatim}

    \begin{doublespace}
        \textbf{\Large{{~\subjectDocument}}}
    \end{doublespace}
\end{center}

\vspace*{\fill}

\begin{flushleft}
    \begin{tabularx}{\linewidth}{@{}>{\bfseries}l@{\hspace{.9em}}X@{}}
        \textbf{Vorgelegt von:} & Andreas Huber <andreas.huber@st.oth-regensburg.de> \\
        \textbf{Matrikelnummer:}& 3370380 \\
        \textbf{Studiengang:}   & Master Informatik (Schwp. Software Engineering) \\
                                & \\
        \textbf{Betreuer:}      & Prof. Dr. techn., Dipl.-Ing. Markus Kucera \\
                                & \\
        \textbf{Abgabedatum:}   & \today \\
    \end{tabularx}
\end{flushleft}
\newpage

% Römische Nummerierung
\pagenumbering{Roman}

% Inhaltsverzeichnis
\tableofcontents

% Abbildungsverzeichnis
\newpage
\phantomsection
\addcontentsline{toc}{section}{Abbildungsverzeichnis}
\renewcommand\refname{Abbildungsverzeichnis}

\listoffigures

% Tabellenverzeichnis
\newpage
\phantomsection
\addcontentsline{toc}{section}{Tabellenverzeichnis}
\renewcommand\refname{Tabellenverzeichnis}

\listoftables

% Abkürzungsverzeichnis
\newpage
\phantomsection
\section*{Abkürzungsverzeichnis}
\addcontentsline{toc}{section}{Abkürzungsverzeichnis}

\begin{acronym}[Bash]
    % \acro{bash}[Bash]{Bourne-again shell}
\end{acronym}

% Arabische Seitennummerierung ab hier
\newpage
\pagenumbering{arabic}

% Einleitung
\section{Einleitung}\label{einleitung}
Visuelle Ähnlichkeitsbestimmung von Bildern spielt in der heutigen Zeit eine
immer wichtigere Rolle. Sie wird neben der Gesichtserkennung beispielsweise auch
für Dublettenerkennung, Anwendung des Urheberrechts-Diensteanbieter-Gesetz,
Qualitätsprüfung und vieles mehr verwendet. Da diese Problematik immer noch ein
offenes Forschungsthema ist und es bisher keine allgemeingültige Lösung gibt,
beschäftigt sich dieses Paper mit einer Übersicht aller aktuell gängigen
Kategorien an \glqq{}Content-Based Image Retrieval\grqq{}-Algorithmen (CBIR).
% TODO: Change!


% Fazit
\newpage
\section{Fazit und Ausblick}\label{fazit-u-ausblick}
Zusammenfassend bietet dieses Paper einen groben Überblick über den aktuellen
Forschungsstand bei der Ähnlichkeitsbestimmung von Bildern. Je nach Größe und
Art des Projekts eignen sich manche Algorithmen mehr und manche weniger. Eine 
allgemeingültige Lösung gibt es bisher nicht.

Der feingranulare inhaltsbasierte Vergleich zwischen Bildern wird auch
sicherlich in Zukunft noch ein wesentlicher Bestandteil der Forschung sein.
Neben den in der Einleitung \ref{einleitung} erläuterten Anwendungsfällen
spielen Image Matching Algorithmen insbesondere in der Medizin eine immer
größere Rolle. Beispielsweise sind computergestützte Diagnosen durch 
Bildvergleichsalgorithmen in der Lage Radiologen bei der Erkennung von
Krankheiten zu helfen \parencite{conclusion-medical}.

% Römische Nummerierung
\newpage
%\pagenumbering{Roman}
%\setcounter{page}{5}

% Literaturliste soll im Inhaltsverzeichnis auftauchen
\newpage
\phantomsection
\addcontentsline{toc}{section}{Quellenverzeichnis}
% Literaturverzeichnis anzeigen
\renewcommand\refname{Quellenverzeichnis}
\printbibliography


% Anhang
\newpage
\phantomsection
\appendix
\section*{Anhang}
\markboth{Anhang}{}
\addcontentsline{toc}{section}{Anhang}

\subsection*{Perceptual Hashing: Implementierung}
\noindent
Schrittweise Generierung eines Perceptual Hashes in Python unter der Verwendung
der OpenCV-Bibliothek.

\begin{lstlisting}[style=Python]
# Perceptual Hash Implementation by Andreas Huber
# Date: 2022-12-10
import numpy as np
import cv2

# Generates Perceptual Hash
def generatePerceptualHash(path, export=False):
    # Step 1: Read Image
    image = cv2.imread(path)
    if (export):
        cv2.imwrite('phash-results/phash-step-1.png', image)

    # Step 2: To grayscale	
    image = cv2.cvtColor(image, cv2.COLOR_BGR2GRAY)
    if (export):
        cv2.imwrite('phash-results/phash-step-2.png', image)

    # Step 3: To 32x32
    image = cv2.resize(image, (32, 32), interpolation = cv2.INTER_NEAREST)
    if (export):
        cv2.imwrite('phash-results/phash-step-3.png', image)

    # Step 4a: Discrete Cosinus Transform: Rows
    image = np.float32(image) / 255.0
    rows = cv2.dct(image, cv2.DCT_ROWS)
    
    # Step 4b: Discrete Cosinus Transform: Columns
    rows = cv2.rotate(rows, cv2.ROTATE_90_COUNTERCLOCKWISE)
    rows = cv2.flip(rows, 0)
    cols = cv2.dct(rows, cv2.DCT_ROWS)
    image = cv2.flip(cols, 0)
    image = cv2.rotate(image, cv2.ROTATE_90_CLOCKWISE)

    if (export):
        cv2.imwrite('phash-results/phash-step-4.png', image * 255)

    # Step 5: Crop to 8x8
    image = image[0:8, 0:8]

    if (export):
        cv2.imwrite('phash-results/phash-step-5.png', image * 255)

    # Step 6a: Calculate median
    imageValues = image.flatten()
    median = np.median(imageValues)

    # Step 6b: intensity < median ? black : white
    for i in range(8):
        for j in range(8):
            if (image[i][j] < median):
                image[i][j] = 0.0
            else:
                image[i][j] = 1.0

    if (export):
        cv2.imshow('resultHash', image)
        cv2.imwrite('phash-results/phash-step-6.png', image * 255)

    # Build string binary perceptual hash
    hash = image.flatten()
    hash = map(int, hash)

    if (export):
        cv2.waitKey(0)
        cv2.destroyAllWindows()

    return ''.join(map(str, hash))
\end{lstlisting}

\end{document}